\section{Módulos}
\subsection{Módulos y morfismos}

Sea $R$ un anillo un $R$-módulo a izquierda consiste de un grupo abeliano $M$ junto
con una operación $\cdot : R \times M \to M$, $(r, m) \mapsto r \cdot m : \quad \forall m$, $n \in M$
y $\forall r$, $s \in R$ se satisfacen las siguientes propiedades:
\begin{enumerate}
    \item $1_R \cdot m = m$.
    \item $(rs) \cdot m = r \cdot ( s \cdot m)$.
    \item $(r + s) \cdot m = r \cdot m + s \cdot m$.
    \item $r \cdot (m + n) = r \cdot m + r \cdot n$.
\end{enumerate}

\begin{note}
    Dado un $R$-módulo a izquierda $M$ tenemos: \begin{enumerate}
        \item $0_R \cdot m = 0_M \quad \forall m \in M$.
        \item $(-r) \cdot m = - (r \cdot m) \quad \forall m \in M$.
        \item $End_{\Z}(M) = \{ f : M \to M : f \text{ es morfismo de grupos} \}$ es un anillo con las operaciones $(f+g)(m) = f(m) + g(m)$ y $(f \cdot g)(m) = f(g(m))$. El neutro para la suma es la función nula y para $\cdot$ es la identidad.
    \end{enumerate}
\end{note}

La operación $\cdot : R \times M \to M$ induce un morfismo de anillos $\rho : R \to End_{\Z}(M)$, $\rho(r)(m) = r \cdot m$.
Recíprocamente, dados un anillo $R$ y un grupo abeliano $M$, un morfismo de anillos $\rho : R \to End_{\Z}$ determina una estructura de $R$-módulo (a izquierda) en $M$, pues
$r \cdot m = \rho(r)(m) \quad \forall r \in R$, $m \in M$. Queda como ejercicio verificar que las propiedades de un $R$-módulo a izquierda se cumplen.

Luego dado un anillo $R$ y un grupo abeliano $M$, dar una estructura de $R$-módulo a izquierda en $M$ es equivalente a dar un morfismo de anillos $\rho : R \to End_{\Z}(M)$.

\begin{definition}[Anulador]
    El anulador de un $R$-módulo $M$ es el ideal $Ann_R(M) = ker(\rho) \triangleleft R$.

    $Ann_R(M) = \{ r \in R : r \cdot m = 0_M \quad \forall m \in M \}$.

    Decimos que $M$ es fiel si $Ann_R(M) = \{0_R\}$ i.e $\rho$ es monomorfismo.
\end{definition}

\begin{definition}[Submódulo]
    Un submódulo de $R$-módulo $M$ es un subgrupo $N \subseteq M : R \cdot N \subseteq N$.
\end{definition}

\begin{definition}[$R$-módulo a derecha]
    Un $R$-módulo a derecha es un grupo abeliano $M$ junto con una operación $\cdot : M \times R \to M$, $(m$, $r) \mapsto m \cdot r$ que verifica: \begin{enumerate}
        \item $m \cdot 1_R = m$.
        \item $m \cdot (rs) = (m \cdot r) \cdot s$.
        \item $m \cdot (r + s) = m \cdot r + m \cdot s$.
        \item $(m + n) \cdot r = m \cdot r + n \cdot r$.
    \end{enumerate}
    $\forall m$, $n \in M$ y $\forall r$, $s \in R$
\end{definition}

Notemos que esto es equivalente a que $M$ tenga una estructura de $R^{op}$-módulo a izquierda:

$\tilde{\cdot} \, : R^{op} \times M \to M$, $(r$, $m) \mapsto m \, \tilde{\cdot} \, r$.
\begin{enumerate}
    \item $1_R \, \tilde{\cdot} \, m = m$.
    \item $(rs) \, \tilde{\cdot} \, m = sr \, \tilde{\cdot} \, m = (m \cdot s) \cdot r = r \, \tilde{\cdot} \, (m \cdot s) = r \, \tilde{\cdot} \, (s \, \tilde{\cdot} \, m)$.
    \item $(r + s) \, \tilde{\cdot} \, m = m \cdot (r + s) = m \cdot r + m \cdot s = r \, \tilde{\cdot} \, m + s \, \tilde{\cdot} \, m$.
    \item $ r \, \tilde{\cdot} \, (m + n) = (m + n) \cdot r = m \cdot r + n \cdot r = r \, \tilde{\cdot} \, m + r \, \tilde{\cdot} \, n$.
\end{enumerate}

Recíprocamente, si $M$ es un $R^{op}$-módulo a izquierda $\cdot : R^{op} \times M \to M$, define $* : M \times R \to M$ por $m * r := r \cdot m$. Verificar que cumple las propiedades de $R$ módulo a derecha.

\begin{definition}
    Sean $R$ un anillo, $M$ y $N$ dos $R$-submódulos y $f : M \to N$ una función. Decimos que $f$ es un morfismo de $R$-módulos si $f$ es morfismo de grupos (morfismo $R$-lineal) y $f(r \cdot m) = r \cdot f(m) \quad \forall r \in R \quad \forall m \in M$.
\end{definition}

Notaremos \begin{enumerate}
    \item $hom_R(M$, $N) = \{ f : M \to N : f \text{ es morfismo} \}$. Es grupo abeliano con la suma puntual.
    \item $End_R(M) = hom_R(M$, $M)$.
    \item $Aut_R(M) = \{ f \in End_R(M) : f \text{ es biyectivo} \}$.
\end{enumerate}

\begin{note}
    Si $P$ es otro $R$-módulo, la composición de una aplicación $\circ : hom_R(P$, $N) \times hom_R(M$, $P) \to hom_R(M$, $N)$. Es decir, la composición de morfismos $R$-lineales es un morfismo $R$-lineal.
    Además, la aplicación $\circ$ es bilineal, es decir, $\forall f$, $g \in hom_R(M$, $N)$ y $\forall h$, $k \in hom_R(P$, $M)$ se cumple que:
    \begin{align*}
        (f + g) \circ h & = f \circ h + g \circ h \\
        f \circ (h + k) & = f \circ h + f \circ k \\
    \end{align*}
\end{note}

En particular, el grupo $End_R(N)$ junto con la composición es un anillo y $hom_R(M$, $N)$ es un $End_R(N)$-módulo a izquierda y un $End_R(M)$-módulo a derecha.
$f \in hom_R(M$, $N)$, $g \in End_R(N)$, $h \in End_R(M)$ entonces $g \circ f = g \cdot f$, $f \circ h = f \cdot h$.

Si $a \in Z(R)$, la aplicación $L_a : N \to N$ dada por $L_a(n) = a \cdot n$ es $R$-lineal y $L: Z(R) \to End_R(N)$ es morfismo de anillos.

Esto nos permite darle a $hom_R(M$, $N)$ una estructura de $Z(R)$-módulo.

\begin{definition}
    Sea $M$ un grupo abeliano y $R$, $S$ anillos. Decimos que $M$ es un $(R$, $S)$-bimódulo si es un $R$-módulo a izquierda y un $S$-módulo a derecha y las estructuras son compatibles.
    Es decir, $\forall r \in R \, \forall s \in S \, \forall m \in M$, $(r \cdot m) \cdot s = r \cdot (m \cdot s)$.
\end{definition}